\Needspace{15\baselineskip}
\subsection{The Anointing of the Sick}

\begin{massbox}{Formal account: the Church's prayer in oil}
Anointing of the Sick is the sacrament in which a bishop or presbyter anoints a seriously ill baptized member with blessed oil while pronouncing the Church's approved prayer, so that the suffering and glorified Christ may confer the Holy Spirit's strength, forgive sins when needed and receivable, relieve the sick person, restore bodily health when conducive to salvation, and prepare the whole person for resurrection and glory. Oil and anointing are the sensible matter; the priestly prayer is the determining form; Christ is principal physician and the ordained minister his living ecclesial instrument (Jas. 5:14--15; Trent, sess.~XIV, \emph{De extrema unctione}, chs.~1--3; CCC 1499, 1520--1523).
\end{massbox}

\subsubsection{Sickness inside creation and covenant}

Sickness is not a positive substance created by God, nor is every illness the direct penalty of a particular personal sin. It is a privation within a good but mortal organism: the failure of powers and organs to possess or exercise the order due to embodied life. Scripture refuses the friends' claim that Job's suffering proves hidden wickedness and records Christ's rejection of a simple guilt calculus for the man born blind (Job 1--2; 42:7--9; John 9:1--3). Yet sickness belongs to the wider dominion of corruption that entered human history with sin, and it exposes a person's dependence, finitude, disordered attachments, fear, and eventual death (Gen. 3:16--24; Rom. 5:12; 8:18--23).

The Old Covenant holds together natural care, prayer, and divine lordship. Oil softens wounds; priests discern conditions affecting worship and communal life; prophets intercede; Hezekiah both prays and receives a poultice; Tobit is healed through a creature used under angelic direction (Isa. 1:6; Lev. 13--14; 2~Kgs. 20:1--11; Tob. 11:7--15). Sirach explicitly commands honor for the physician and medicine because the Lord created both, while commanding prayer and moral rectitude (Sir. 38:1--15). Sacramental anointing therefore need not compete with medicine: nature's arts and the sacrament operate in different but ordered modes under one providence.

Oil bears a dense covenantal history. It nourishes and gives light; it consecrates priests, kings, sanctuary, and vessels; it marks royal and prophetic mission; it signifies gladness, abundance, strengthening, and the Spirit (Exod. 30:22--33; 1~Sam. 16:1--13; Ps. 44:8, Vulgate; modern 45:8; Isa. 61:1--3). The sick person's oil is not chrism and does not confer another character, but it belongs to the same providential grammar: the Anointed One uses created oil to communicate invisible strength to an embodied member.

\subsubsection{Christ heals, sends, suffers, and rises}

The Gospels never reduce Christ's cures to demonstrations of power. He is moved with compassion, touches what is socially untouchable, restores persons to households and worship, and joins visible healing to the authority to forgive (Mark 1:40--45; 2:1--12; 5:21--43). The healings announce the kingdom and disclose who Jesus is, yet they remain provisional signs: every healed person later dies. Their final referent is the Pasch, in which Christ does not evade suffering but bears it, defeats death from within, and rises bodily beyond corruption (Isa. 53:4--5; Matt. 8:14--17; 1~Cor. 15:20--28, 42--57).

Mark 6:7--13 gives the apostolic mission a bodily form: the Twelve preach repentance, expel demons, and anoint many sick persons with oil and heal them. The Catholic tradition does not isolate the oil from the sending. The men act under Christ's authority; word, conversion, victory over evil, and bodily healing belong to one sign of the approaching reign. Trent says the sacrament was ``insinuated'' in Mark and promulgated by James, thereby distinguishing a Gospel anticipation from the apostolic directive that displays the ecclesial rite (sess.~XIV, \emph{De extrema unctione}, ch.~1 and can.~1).

James 5:13--18 is the controlling New Testament dossier. Suffering calls forth personal prayer; cheerfulness, psalmody; sickness, the summoning of the Church's presbyters. They pray \emph{over} the sick person and anoint with oil ``in the name of the Lord.'' The prayer of faith will save the sick, the Lord will raise the person, and sins, if present, will be forgiven. The verbs resist reduction. ``Save'' can embrace bodily rescue and eschatological salvation; ``raise'' can name recovery while opening toward resurrection; conditional forgiveness keeps sickness and personal guilt from being simply identified. The named agents are likewise ordered: the sick summon, the presbyters pray and anoint, the Church believes, but the Lord raises.

\begin{fourstable}{Scriptural phrase}{What it establishes}{Sacramental development}{Limit}
``Presbyters of the Church'' & The act is ecclesial and entrusted to ordained presbyters, not simply to any older or devout person. & Bishops and presbyters alone validly minister the sacrament. & Every Christian may pray for healing and may use blessed oil devotionally where approved; that does not make every anointing sacramental.\\
``Anointing him with oil'' & A material, bodily act belongs to the apostolic directive. & Blessed oil applied by anointing is the sensible bearer of Christ's healing act. & Oil is neither magic nor a pharmaceutical substitute; sacramental efficacy comes from Christ's institution.\\
``Prayer of faith'' & The Church's supplication is intrinsic, not an optional commentary on oil. & The priestly prayer determines the anointing toward saving, raising, forgiveness, and relief. & Efficacy does not depend on the minister's emotional confidence or holiness.\\
``If he be in sins'' & Forgiveness truly belongs among the effects while remaining conditional in James's syntax. & The sacrament forgives sins when the sick cannot confess and places no obstacle; it always gives its proper strengthening to the disposed recipient. & It is not the ordinary replacement for Penance when conscious sacramental confession can be made.\\
\end{fourstable}

\subsubsection{Patristic and early liturgical reception}

Early Christian oil prayers are important but must be classified with care. The church order traditionally called the \emph{Apostolic Tradition}, ch.~5, contains a blessing over offered oil asking health for those who use it. Its composite text and attribution to Hippolytus remain debated, and the prayer does not by itself furnish a complete sacramental ordo. It is best treated as early liturgical evidence that consecrated oil, ecclesial prayer, bodily healing, and spiritual benefit were held together, not as proof of every later ritual detail.

Origen, preserved here in Rufinus's Latin translation, directly joins confession, priestly medicine, James's presbyters, anointing, prayer, and forgiveness while enumerating the ways sins are remitted (\emph{Homilies on Leviticus} 2.4). Because his immediate topic is penitential remission rather than a systematic distinction between Penance and Anointing, the passage shows an early integrated healing practice without settling every later scholastic allocation of effects.

Pope St.~Innocent I's Letter 25 to Decentius, 8.11 (AD 416; \emph{Si instituta ecclesiastica}; DS 216) is the first explicit magisterial discussion. Answering a question about James, Innocent refuses to exclude bishops from the ``presbyters'' who anoint and speaks of oil consecrated by the bishop and usable for the sick. His allowance for the faithful to use blessed oil cannot be turned into proof that lay persons sacramentally minister Anointing. The text addresses both ecclesial blessing and devotional application in a disciplinary setting; the Church's constant doctrinal judgment identifies bishop or presbyter as sacramental minister (CDF, \emph{Note on the Minister of the Sacrament of the Anointing of the Sick}, 2005, commentary 5--10).

St.~John Chrysostom adduces James 5:14--15 precisely to show the post-baptismal healing entrusted to priests (\emph{On the Priesthood} III.6). The natural parents generate children for this life, he argues, while priests minister birth and reconciliation ordered to the life to come. His purpose is sacerdotal responsibility, not an exhaustive commentary on effects, but he directly protects two constitutive points: sickness after Baptism still needs grace, and James's ecclesial presbyters act in an ordained capacity.

St.~Caesarius of Arles, \emph{Sermon} 13.3, urges the sick not to seek charms and diviners but to approach the Church, receive Christ's Body and Blood, and be anointed with blessed oil according to James. The polemical setting clarifies sacramental causality: oil is not one sacred technique among occult rivals; it belongs to faith, ecclesial prayer, conversion, and Eucharistic communion. The Venerable Bede's \emph{Commentary on the Seven Catholic Epistles}, on James 5:14--15, is direct exegesis of the presbyteral anointing and a further Western witness to its continuity.

The ancient liturgies add what a chain of treatises cannot. East and West bless oil, invoke the Spirit, name sins and weakness, read the Gospel, and surround the sick with ecclesial intercession. Eastern traditions often retain a solemn service with several priests; CCEO 737 \S2 commends preservation of that custom where it exists, but one priest remains sufficient for validity. The older Roman rite distributed anointings across the bodily senses with the repeated petition that the Lord forgive sins committed through each; the current Roman rite ordinarily anoints forehead and hands with a prayer for the Spirit's grace, saving, and raising. Ritual plurality manifests one mystery rather than several rival sacraments.

\begin{studybox}{Direct witnesses and illuminating witnesses}
Mark 6 and James 5 are the scriptural foundation; Chrysostom and Bede directly receive James; Innocent directly rules on its ecclesial administration; Caesarius directly exhorts the sick to the Church's blessed oil. The \emph{Apostolic Tradition} prayer is related early liturgical evidence with a disputed textual history. Origen directly cites James within a penitential dossier, but his passage should not be forced into all later distinctions between the two healing sacraments.
\end{studybox}

\subsubsection{Matter, form, and ritual unity}

The remote matter is blessed oil; the proximate matter is its sacramental application by anointing the sick person. Olive oil is the traditional matter and, for Aquinas, especially fitting because it penetrates, softens, and signifies a gentle, hope-sustaining cure (Supplement, q.~29, a.~4). In the present Roman rite, Paul VI authoritatively permits olive oil or, according to circumstances, another oil derived from plants; it must be duly blessed. The Latin bishop ordinarily blesses the Oil of the Sick, while a presbyter may do so in necessity during the celebration. In Eastern Catholic discipline the administering priest ordinarily blesses the oil unless the particular law of the Church \emph{sui iuris} provides otherwise (Paul VI, \emph{Sacram unctionem infirmorum}; CIC 999; CCEO 741).

The form is a priestly prayer that determines the natural sign of oil toward Christ's saving effects. James gives prayer such prominence that the Supplement calls deprecation the sacrament's formal mode (q.~29, aa.~7--9). No one vernacular paraphrase or Latin minimum should be universalized across Catholic rites: validity requires the form prescribed by the competent Church and a signifying unity of word and anointing. In the current Roman rite the integral sign ordinarily includes anointing forehead and hands with the promulgated formula; in necessity one anointing on the forehead or another suitable part with the whole formula suffices (CIC 998, 1000; Paul VI, \emph{Sacram unctionem infirmorum}).

Several anointings do not make several sacraments. A complete action can possess real parts ordered to one signified cure, just as a house is one whole composed of parts. The medieval Roman anointing of eyes, ears, nostrils, mouth, hands, and feet dramatized the powers through which human action proceeds; newer and Eastern rites realize the sacramental unity differently. Aquinas's specific claim that each senses-anointing was essential described the rite and theology available to him; it is not a limit on the Church's authority to determine which elements preserve the divinely instituted substance (Supplement, q.~29, a.~2; Paul VI, \emph{Sacram unctionem infirmorum}).

\begin{threestable}{Causal order}{Creaturely instrument}{What the order excludes}
Christ the suffering and glorified Lord & His humanity is the conjoined instrument; his Passion merits and causes sacramental grace; the risen Christ is principal physician. & Oil does not act as magic, and the priest does not heal by private charisma.\\
The Church praying in her priest & The ordained minister applies oil and form \emph{in persona Ecclesiae}; the Church's prayer surrounds the act. & The priest's sin does not nullify Christ's action; the prayer is more than a merely private petition.\\
The sick member receiving & Faith, contrition, intention at least habitual or implicit, and nonresistance dispose the person for fruit. & Disposition does not turn grace into an achievement; absence of bodily cure does not show sacramental failure.\\
\end{threestable}

\subsubsection{Sacramental operation and the hierarchy of ends}

A textual note prevents a false precision. Aquinas's authentic \emph{Summa contra Gentiles} IV.73 gives a compact treatment of Anointing. The detailed Supplement, qq.~29--33, was compiled after his death from his earlier \emph{Commentary on the Sentences}; it is a major Thomistic witness but not the completed autograph continuation of ST III, which stops within the Penance material. Its medieval discipline must therefore be read with the later magisterial determination of the rite, while its causal and teleological distinctions remain illuminating.

The Supplement distinguishes spiritual death from spiritual debility. Baptism is regeneration; Penance is resurrection after mortal sin; Anointing is healing that normally presupposes spiritual life and attacks weakness, proneness toward evil, difficulty in good, fear, and other ``remnants'' left by sin and intensified by illness (q.~30, a.~1; \emph{SCG} IV.73). ``Remnants'' do not mean a surviving guilt that absolution failed to remit. They include disordered dispositions, temporal debt, diminished vigor, and unfinished purification that impede ready action toward glory.

The causal series and the hierarchy of ends must be read together without being collapsed. \emph{Sacramentum}, \emph{res et sacramentum}, and \emph{res tantum} distinguish the sign, its intermediate inward reality, and the grace effected. Primary and secondary ends state why those effects are conferred and which goods govern the others. Bodily recovery is therefore neither the sacrament's defining \emph{res} nor the test of its success.

\begin{fourstable}{Ordered level}{Sacramental operation}{Teleological status}{Controlling locus}
Outward sign & Blessed oil is sacramentally applied to the sick while the priest pronounces the Church's approved prayer of faith. Word and anointing form one sensible ecclesial instrument. & Christ the suffering and risen physician acts through this sign; neither oil's natural properties nor the minister's private charisma causes sacramental grace. & Jas. 5:14--15; Trent XIV, chs.~1--3; Supplement, q.~29, aa.~4, 7--9.\\
Thomistic \emph{res et sacramentum} & The outward anointing signifies and effects an inward devotion or ``spiritual anointing,'' a filial readiness by which the sick member is strengthened under the Spirit. & This intermediate sign-and-reality is not an indelible character; it orders the recipient toward the sacrament's healing gifts and toward glory. & Supplement, q.~30, a.~3 ad 3; \emph{SCG} IV.73.\\
\emph{Res tantum} & In a disposed recipient, sacramental grace heals spiritual debility and sin's remnants, strengthens hope and resistance, unites suffering to Christ, remits sins when needed and receivable, and prepares the person for passage. & These spiritual goods constitute the proper grace effected. Their reality does not depend on felt consolation, cognitive vigor, or physical improvement. & Supplement, q.~30, aa.~1, 3; Trent XIV, ch.~2; CCC 1520--1523.\\
Primary proper / proximate end & Spiritual strengthening and healing: grace fortifies the sick, removes impediments and remnants, forgives guilt when necessary and no obstacle is placed, and prepares the whole person for death and glory. & This primary complex governs the sacrament whether illness ends in recovery or death. Forgiveness remains conditional in James and does not displace Penance when confession is possible. & Jas. 5:15; \emph{SCG} IV.73; Supplement, q.~30, a.~1; Trent XIV, ch.~2; CCC 1520, 1523.\\
Subordinate secondary ends & Relief of bodily suffering and restoration of bodily health may be given when, and only when, they conduce to the recipient's spiritual salvation; ecclesial consolation and offered suffering also serve the primary healing. & Bodily recovery is contingent and secondary, never promised. Its absence proves neither invalidity, unfruitfulness, deficient faith, nor ministerial failure. Ordinary medicine remains a genuine providential good. & Supplement, q.~30, a.~2; Trent XIV, ch.~2; CCC 1519--1522; Sir. 38:1--15.\\
Common ultimate end & Resurrection of the body, perfected charity, and the vision of God; near death this orientation converges with Penance and reaches its sacramental summit in Viaticum. & Anointing's proper and secondary ends serve the same beatitude and divine glory as every sacrament. Temporary recovery remains penultimate; incorruptible life is the final cure. & ST III, q.~60, a.~3; q.~65, a.~3; CCC 1523--1525; 1~Cor. 15:42--57.\\
\end{fourstable}

\begin{fourstable}{Effect}{Thomistic articulation}{Magisterial articulation}{Guardrail}
Strength and grace & The sacrament principally heals spiritual debility and confers grace for the acts of life and glory (Supplement, q.~30, a.~1). & The Holy Spirit gives strength, peace, courage, and resistance to discouragement and the final enemy (Trent XIV, ch.~2; CCC 1520). & Grace need not remove pain, cognitive decline, or fear as felt emotion; it strengthens the person within them.\\
Forgiveness & Grace removes the remnants of sin and, consequently, guilt if it encounters guilt and no obstacle in the recipient. & Sins are forgiven when this is needed, especially when the sick person cannot obtain Penance (Jas. 5:15; CCC 1520). & A conscious person able to confess grave sin should ordinarily receive Penance; Anointing is no strategy for avoiding confession or conversion.\\
Bodily relief & Bodily healing is a real secondary effect when required for spiritual health; divine wisdom acts toward the principal end (Supplement, q.~30, a.~2). & The sacrament may restore health if conducive to salvation (Trent XIV, ch.~2; CCC 1520). & No minister can promise cure, infer deficient faith from continued illness, or oppose ordinary medical care.\\
Preparation for passage & Because the sacrament removes impediments and strengthens, it prepares the sick for glory; it may be repeated because its remedial effect is not a character (\emph{SCG} IV.73; Supplement, q.~30, a.~3; q.~33). & It completes holy anointings marking Christian life and fortifies the end of the earthly pilgrimage (CCC 1523). & It is not reserved to the final minutes; timely reception permits conscious faith and pastoral care.\\
\end{fourstable}

The Supplement assigns no character because Anointing deputes no one to a permanent cultic office. Its \emph{res et sacramentum} is instead a certain inward devotion or ``spiritual anointing'' produced and signified by the outward rite (q.~30, a.~3 ad 3). This school allocation is not the only later account, but it captures a vital truth: the sacrament does not merely alter an external status. The Spirit interiorly conforms the sick member to Christ's willing, filial endurance and thereby points toward further effects of forgiveness, relief, and glory.

The union with Christ's Passion is neither a cult of pain nor a claim that suffering adds deficient merit to Calvary. Christ's oblation is all-sufficient. Because the baptized are real members of his Body, however, grace can make their received suffering an act of charity within his one self-offering. Colossians 1:24 concerns participation in the afflictions of Christ's ecclesial Body, not completion of an insufficient redemption. The sick person, often forced into apparent passivity, is thus not expelled from mission: prayer, patient endurance, truthful lament, and offered dependence can benefit the whole Church (CCC 1499, 1521--1522).

\subsubsection{Minister, subject, repetition, and limits}

Every bishop or presbyter, and only a bishop or presbyter, validly administers Anointing. James names the Church's presbyters; Trent defines the proper minister as priest; Latin and Eastern canons agree. The CDF states that this teaching is to be held definitively: a deacon or lay person cannot be deputed as an ``extraordinary minister,'' and an attempted sacramental anointing by either is simulation, not a valid but illicit sacrament (Trent, sess.~XIV, can.~4; CIC 1003 \S1; CCEO 739 \S1; CDF, \emph{Note}, 2005).

This validity rule differs from the blessing and devotional use of oil. Lay persons may pray with the sick, touch them appropriately, and use approved sacramentals; deacons may preach, bring Communion, and lead prayers. None thereby pronounces the sacramental form as minister. Nor does a priest need a Penance-like faculty for validity. Pastoral office and permission govern ordinary liceity, but in necessity any priest may and should administer according to law; a Latin presbyter may bless the oil in the actual celebration when necessary (CIC 999, 1003; CCEO 739--741).

The subject is a baptized living person seriously endangered by illness or old age. The current Latin threshold is reached when a faithful person who has attained the use of reason \emph{begins} to be in danger from sickness or age; one need not await imminent death. Eastern law speaks of the gravely ill and presumes the desire of an unconscious faithful person in danger or according to priestly judgment (CIC 1004--1006; CCEO 737--740). A healthy person facing execution or dangerous travel is not a subject, because the sign is bodily medicine applied to bodily illness. Pregnancy, surgery, disability, mental illness, or old age are not automatic titles in abstraction; the actual serious impairment or danger is pastorally judged.

\begin{studybox}{Practical sacramental grammar}
\textbf{Doubt favors care.} Latin law directs anointing when there is doubt whether the person attained reason, is dangerously ill, or is still alive; no sacrament is given after death is certain (CIC 1005). \textbf{Implicit desire can suffice.} One who requested the sacrament at least implicitly while capable is to receive it when unconscious (CIC 1006). \textbf{Obstinacy remains an obstacle.} It is not conferred on one obstinately persevering in manifest grave sin; this is not a judgment that illness proves sin, but respect for a sacrament whose inward direction is conversion (CIC 1007). \textbf{Repetition is medicinal.} It may be repeated after recovery and a new grave illness, or when danger becomes graver during the same illness (CIC 1004 \S2; Supplement, q.~33). \textbf{Approved rite governs.} Matter, blessing, number and place of anointings, and formula follow the competent Latin or Eastern liturgical books; emergency abbreviations are those books' provisions, not private invention.
\end{studybox}

\subsubsection{The last sacraments and the glory of the body}

When circumstances allow, the sacramental order near death is Penance, Anointing, and Viaticum. Penance reconciles the conscious sinner; Anointing strengthens and completes healing amid serious illness; the Eucharist gives the pilgrim Christ himself as food for passage. ``Last rites'' is a pastoral complex, not another sacrament, and Anointing is not the Eucharist. Viaticum is preeminently the last sacrament because the Bread of Life is pledge of resurrection and seed of glory (John 6:39--58; CCC 1524--1525).

The other sacraments converge here. Baptism began conformity to death and resurrection; Confirmation anointed for witness; Penance restored lost grace; Orders gives the priestly instrument who now speaks the Church's prayer; Matrimony and family life often furnish the communion that carries the sick member; the Eucharist brings the Head substantially to the member. Anointing adds no character because it does not initiate a new ecclesial function. It strengthens the already initiated person to bring baptismal configuration through sickness toward its eschatological term.

James's ``the Lord will raise him up'' finally exceeds every temporary recovery. A healed body can again serve, love, suffer, and die; the risen body cannot die again. The sacrament therefore honors the present body without absolutizing present biological survival. It asks relief with confidence, receives medicine without superstition, accepts palliative and ordinary care without treating death as a technical failure, and hopes for the transformation in which corruption puts on incorruption (1~Cor. 15:42--57).

\begin{massbox}{Oil at the boundary of the ages}
The same creature that soothed wounds, fed lamps, and consecrated Israel's servants is taken into Christ's final medicine. The Church places it on a mortal body and prays to the suffering \emph{and glorified} Lord. Thus the sign holds together what fear divides: sin may be forgiven, weakness strengthened, the body relieved, suffering offered, death approached, and resurrection expected. The sacrament does not call death good. It declares that even at death's frontier the baptized body remains Christ's member and is destined not for disposal but for glory.
\end{massbox}
